% unit: u23
% sources: sources/095-eigenvalues-and-eigenvectors.txt, sources/041-lecture-21-eigenvalues-and-eigenvectors.txt, sources/194-mit18-06scf11-ses2-8sum-pdf-download-file.txt, sources/042-problem-solving-eigenvalues-and-eigenvectors.txt, sources/195-mit18-06scf11-ses2-8prob-pdf-download-file.txt, sources/196-mit18-06scf11-ses2-8sol-pdf-download-file.txt
\section{Eigenvalues and eigenvectors}\label{sec:u23}

Up to now a matrix has mostly been the coefficient matrix of a static problem
$A\vb{x}=\vb{b}$. From here on we also think of a square matrix as an operation that is
applied again and again: $A\vb{u}$, $A^{2}\vb{u}$, $A^{3}\vb{u}$, or continuously in time.
For that, the important vectors are the few whose direction $A$ does not change. They and
their stretching factors, the eigenvectors and eigenvalues, occupy most of the rest of the
course. This section says what they are and how to find them; the next sections use them.

\subsection{Vectors that keep their direction}

A matrix acts like a function: in goes $\vb{x}$, out comes $A\vb{x}$. For most vectors the
output points in a new direction. The eigenvectors are the exceptions.

\begin{definition}[Eigenvalue and eigenvector]\label{def:u23-eig}
Let $A$ be $n\times n$. A number $\lambda$ is an \emph{eigenvalue} of $A$ if there is a
\emph{nonzero} vector $\vb{x}$ with
\[
  A\vb{x}=\lambda\vb{x}.
\]
Such an $\vb{x}$ is an \emph{eigenvector} for $\lambda$. The eigenvectors for $\lambda$,
together with $\vb{0}$, form the \emph{eigenspace} $\Nul{A-\lambda I}$.
\end{definition}

``Same direction'' allows $\lambda$ to be negative (the vector reverses), zero, or even
complex. The vector $\vb{0}$ is excluded because $A\vb{0}=\lambda\vb{0}$ for every
$\lambda$ and would tell us nothing. Any nonzero multiple of an eigenvector is again an
eigenvector, so ``the'' eigenvector really means a basis for the eigenspace.

\emph{Eigenvalue zero.} $A\vb{x}=0\vb{x}$ says $\vb{x}$ is a nonzero vector in the nullspace.
So $\lambda=0$ is an eigenvalue exactly when $A$ is singular, and its eigenvectors are the
nonzero vectors of $\Nul{A}$. Among eigenvalues, zero is no longer special.

Before any formulas, two matrices whose eigenvectors we can see from what they do.

\begin{example}[A projection matrix]\label{ex:u23-projection}
Let $P$ project $\R^{3}$ onto the plane $x+y+z=0$. By the formula
$P=I-\vb{n}\vb{n}\T/\vb{n}\T\vb{n}$ with normal $\vb{n}=(1,1,1)$,
\[
  P=\frac13\begin{bmatrix} 2&-1&-1\\ -1&2&-1\\ -1&-1&2\end{bmatrix}.
\]
A vector already in the plane is not moved: $P\vb{x}=\vb{x}$, so every nonzero vector of the
plane is an eigenvector with $\lambda=1$, for example $(1,-1,0)$ and $(1,0,-1)$. A vector
along the normal is projected to zero: $P(1,1,1)=\vb{0}$, so $\vb{n}$ is an eigenvector with
$\lambda=0$. The eigenvalues are $1,1,0$, and the three eigenvectors are independent. A
typical vector such as $\vb{b}=(1,0,0)$ is \emph{not} an eigenvector:
$P\vb{b}=\tfrac13(2,-1,-1)$ points in another direction.
\end{example}

\begin{example}[A permutation matrix]\label{ex:u23-permutation}
$B=\begin{bmatrix} 0&1\\ 1&0\end{bmatrix}$ swaps the two components:
$B(x_1,x_2)=(x_2,x_1)$. The vector $(1,1)$ is unchanged, so it is an eigenvector with
$\lambda=1$. The vector $(1,-1)$ becomes $(-1,1)$, its negative, so it is an eigenvector
with $\lambda=-1$. A matrix with no negative entries can still have a negative eigenvalue.
\end{example}

\subsection{The characteristic equation}

The equation $A\vb{x}=\lambda\vb{x}$ has two unknowns multiplied together, so elimination
cannot be applied to it directly. The trick is to find $\lambda$ first.

\begin{theorem}[Characteristic equation]\label{thm:u23-chareq}
$\lambda$ is an eigenvalue of $A$ if and only if
\[
  \det(A-\lambda I)=0 .
\]
The eigenvectors for $\lambda$ are the nonzero vectors in $\Nul{A-\lambda I}$.
\end{theorem}

\begin{proof}
$A\vb{x}=\lambda\vb{x}$ is the same as $(A-\lambda I)\vb{x}=\vb{0}$. A nonzero solution
exists exactly when $A-\lambda I$ is singular, and a square matrix is singular exactly when
its determinant is zero (Section~\ref{sec:u20}).
\end{proof}

By the big formula, $p(\lambda)=\det(A-\lambda I)$ is a polynomial of degree $n$ in
$\lambda$, the \emph{characteristic polynomial}; its leading term comes from the product of
the diagonal entries and is $(-\lambda)^{n}$. Over the complex numbers it factors as
$p(\lambda)=(\lambda_1-\lambda)(\lambda_2-\lambda)\cdots(\lambda_n-\lambda)$, so an
$n\times n$ matrix has exactly $n$ eigenvalues if repeated roots are counted with their
multiplicity. The recipe is:
\begin{enumerate}
  \item compute $\det(A-\lambda I)$ and find its $n$ roots $\lambda_1,\dots,\lambda_n$;
  \item for each $\lambda_i$, find the nullspace of the singular matrix $A-\lambda_i I$ by
  elimination (give the free variable the value $1$).
\end{enumerate}

\begin{example}[A symmetric $2\times 2$ matrix]\label{ex:u23-31}
Let $A=\begin{bmatrix} 3&1\\ 1&3\end{bmatrix}$. Then
\[
  \det(A-\lambda I)=\begin{vmatrix} 3-\lambda & 1\\ 1 & 3-\lambda\end{vmatrix}
  =(3-\lambda)^{2}-1=\lambda^{2}-6\lambda+8=(\lambda-4)(\lambda-2),
\]
so $\lambda_1=4$ and $\lambda_2=2$. For $\lambda_1=4$:
$A-4I=\left[\begin{smallmatrix} -1&1\\ 1&-1\end{smallmatrix}\right]$, whose nullspace is the
line through $\vb{x}_1=(1,1)$. For $\lambda_2=2$:
$A-2I=\left[\begin{smallmatrix} 1&1\\ 1&1\end{smallmatrix}\right]$, whose nullspace is the
line through $\vb{x}_2=(-1,1)$. Check: $A(1,1)=(4,4)$ and $A(-1,1)=(-2,2)$.
Figure~\ref{fig:u23-stretch} shows the two eigenvector lines.
\end{example}

\begin{figure}[ht]\centering
\begin{tikzpicture}[scale=1.0,>=Stealth]
  \draw[->,gray] (-3,0) -- (5,0) node[right] {$x_1$};
  \draw[->,gray] (0,-0.5) -- (0,4.8) node[above] {$x_2$};
  \draw[blue!50,dashed] (-0.5,-0.5) -- (4.6,4.6);
  \draw[red!50,dashed] (0.5,-0.5) -- (-2.9,2.9);
  \draw[->,very thick,blue!60] (0,0) -- (4,4) node[above left] {$A\vb{x}_1=4\vb{x}_1$};
  \draw[->,very thick,blue!90!black] (0,0) -- (1,1) node[above left] {$\vb{x}_1$};
  \draw[->,very thick,red!60] (0,0) -- (-2,2) node[left] {$A\vb{x}_2=2\vb{x}_2$};
  \draw[->,very thick,red!80!black] (0,0) -- (-1,1) node[below left] {$\vb{x}_2$};
  \draw[->,thick,black!70] (0,0) -- (1,0) node[below] {$\vb{v}$};
  \draw[->,thick,black!70,dashed] (0,0) -- (3,1) node[right] {$A\vb{v}$};
\end{tikzpicture}
\caption{For $A=\left[\begin{smallmatrix} 3&1\\ 1&3\end{smallmatrix}\right]$, vectors on the
dashed line through $(1,1)$ are stretched by $\lambda_1=4$ and vectors on the line through
$(-1,1)$ by $\lambda_2=2$, without turning. A typical vector such as $\vb{v}=(1,0)$ is turned:
$A\vb{v}=(3,1)$ is not a multiple of $\vb{v}$.}
\label{fig:u23-stretch}
\end{figure}

The matrix in Example~\ref{ex:u23-31} is $B+3I$, with $B$ the permutation of
Example~\ref{ex:u23-permutation}. Its eigenvalues $4,2$ are those of $B$ (namely $1,-1$)
plus $3$, and its eigenvectors are the same lines. This is a general fact.

\begin{proposition}[Shifting by a multiple of $I$]\label{prop:u23-shift}
If $A\vb{x}=\lambda\vb{x}$, then $(A+cI)\vb{x}=(\lambda+c)\vb{x}$. So $A+cI$ has the
eigenvalues of $A$ shifted by $c$, with the same eigenvectors.
\end{proposition}

\begin{proof}
$(A+cI)\vb{x}=A\vb{x}+c\vb{x}=\lambda\vb{x}+c\vb{x}=(\lambda+c)\vb{x}$.
\end{proof}

\begin{trap}
It is tempting to argue: if $A\vb{x}=\lambda\vb{x}$ and $B\vb{x}=\alpha\vb{x}$ then
$(A+B)\vb{x}=(\lambda+\alpha)\vb{x}$, so eigenvalues add. The flaw is that $A$ and $B$
usually have \emph{different} eigenvectors, and then nothing can be concluded. For
$A=\left[\begin{smallmatrix} 0&1\\ 0&0\end{smallmatrix}\right]$ and
$B=\left[\begin{smallmatrix} 0&0\\ 1&0\end{smallmatrix}\right]$ all four eigenvalues are $0$,
yet $A+B$ has eigenvalues $1$ and $-1$, and $AB=\left[\begin{smallmatrix} 1&0\\ 0&0\end{smallmatrix}\right]$
has eigenvalues $1$ and $0$. Eigenvalues are neither additive nor multiplicative; the shift
$B=cI$ works only because every vector is an eigenvector of $cI$.
\end{trap}

\subsection{Trace and determinant}

In Example~\ref{ex:u23-31} the coefficient $6$ in $\lambda^{2}-6\lambda+8$ is $3+3$, the sum
of the diagonal entries, and the constant $8$ is $\det A$. That is no accident.

\begin{theorem}[Sum and product of the eigenvalues]\label{thm:u23-trace}
Let $\lambda_1,\dots,\lambda_n$ be the eigenvalues of $A$, repeated according to
multiplicity. Then
\[
  \lambda_1+\lambda_2+\dots+\lambda_n=a_{11}+a_{22}+\dots+a_{nn}=\tr A,
  \qquad
  \lambda_1\lambda_2\cdots\lambda_n=\det A .
\]
\end{theorem}

\begin{proof}
Compare two expressions for the characteristic polynomial,
\[
  \det(A-\lambda I)=(\lambda_1-\lambda)(\lambda_2-\lambda)\cdots(\lambda_n-\lambda).
\]
\emph{Product.} Set $\lambda=0$: $\det A=\lambda_1\lambda_2\cdots\lambda_n$.

\emph{Sum.} Look at the coefficient of $\lambda^{n-1}$. On the right it is
$(-1)^{n-1}(\lambda_1+\dots+\lambda_n)$. On the left, use the big formula: each term is a
product of $n$ entries of $A-\lambda I$, one from each row and each column. A term that uses
an off-diagonal entry in position $(i,j)$ cannot also use the diagonal entries of row $i$ or
of column $j$, so it contains at most $n-2$ factors with $\lambda$. Hence only the diagonal
product $(a_{11}-\lambda)\cdots(a_{nn}-\lambda)$ reaches $\lambda^{n-1}$, with coefficient
$(-1)^{n-1}(a_{11}+\dots+a_{nn})$. Equate the two coefficients.
\end{proof}

For a $2\times 2$ matrix this pins the characteristic equation down completely:
\[
  \lambda^{2}-(\tr A)\,\lambda+\det A=0 .
\]
Knowing one eigenvalue, the trace gives the other at once. The two facts are also the
quickest check on a computation: add the eigenvalues and compare with the diagonal.

\begin{example}[Two cities: finding eigenvalues by inspection]\label{ex:u23-markov}
Each year $80\%$ of the people in city~1 stay and $20\%$ move to city~2, while $70\%$ of
city~2 stays and $30\%$ moves to city~1. The populations change by
\[
  A=\begin{bmatrix} 0.8&0.3\\ 0.2&0.7\end{bmatrix}.
\]
Try $\vb{x}_1=(600,400)$:
$A\vb{x}_1=(480+120,\;120+280)=(600,400)$. So $\vb{x}_1$ is an eigenvector with
$\lambda_1=1$. The trace is $0.8+0.7=1.5$, so $\lambda_2=1.5-1=0.5$. The determinant
confirms it: $0.56-0.06=0.5=\lambda_1\lambda_2$. For $\lambda_2$,
$A-0.5I=\left[\begin{smallmatrix} 0.3&0.3\\ 0.2&0.2\end{smallmatrix}\right]$ has nullspace
spanned by $\vb{x}_2=(1,-1)$, and indeed $A(1,-1)=(0.5,-0.5)=0.5\,\vb{x}_2$. The two
eigenvectors are the key to what happens to the populations year after year:
Section~\ref{sec:u24} follows them.
\end{example}

\subsection{Complex eigenvalues}

The $90^{\circ}$ rotation $Q=\begin{bmatrix} 0&-1\\ 1&0\end{bmatrix}$ turns every nonzero
vector, so no real vector can be an eigenvector. The trace and determinant agree that
something is strange: $\lambda_1+\lambda_2=0$ but $\lambda_1\lambda_2=1$, impossible for
real numbers. The characteristic equation settles it:
\[
  \det(Q-\lambda I)=\begin{vmatrix} -\lambda & -1\\ 1 & -\lambda\end{vmatrix}=\lambda^{2}+1=0,
  \qquad \lambda_1=i,\quad \lambda_2=-i .
\]
For $\lambda_1=i$, the nullspace of $Q-iI$ contains $\vb{x}_1=(1,-i)$; check:
$Q(1,-i)=(i,1)=i(1,-i)$. Similarly $\vb{x}_2=(1,i)$ goes with $-i$. A perfectly real matrix
can have complex eigenvalues and eigenvectors.

\begin{proposition}[Conjugate pairs]\label{prop:u23-conj}
If $A$ is real and $A\vb{x}=\lambda\vb{x}$, then $A\bar{\vb{x}}=\bar\lambda\bar{\vb{x}}$.
So the nonreal eigenvalues of a real matrix come in conjugate pairs $a\pm ib$, with conjugate
eigenvectors.
\end{proposition}

\begin{proof}
Take complex conjugates of every entry in $A\vb{x}=\lambda\vb{x}$. Since $\bar A=A$, this
gives $A\bar{\vb{x}}=\bar\lambda\bar{\vb{x}}$, and $\bar{\vb{x}}\neq\vb{0}$.
\end{proof}

How far a matrix is from symmetric decides how complex its eigenvalues can be. Symmetric
matrices ($A\T=A$) have real eigenvalues and perpendicular eigenvectors; this is proved in
Section~\ref{sec:u29}, and the examples $B$ and $\left[\begin{smallmatrix} 3&1\\ 1&3\end{smallmatrix}\right]$
above show it. The rotation $Q$ is antisymmetric ($Q\T=-Q$), the opposite extreme, and its
eigenvalues are purely imaginary.

\subsection{Triangular matrices and repeated eigenvalues}

\begin{proposition}[Triangular matrices]\label{prop:u23-triangular}
The eigenvalues of a triangular matrix are its diagonal entries.
\end{proposition}

\begin{proof}
$A-\lambda I$ is triangular with diagonal entries $a_{ii}-\lambda$, and the determinant of a
triangular matrix is the product of its diagonal entries. So
$\det(A-\lambda I)=(a_{11}-\lambda)\cdots(a_{nn}-\lambda)$, with roots $a_{11},\dots,a_{nn}$.
\end{proof}

Triangular matrices show the one real danger in this subject: a repeated eigenvalue can come
with too few eigenvectors.

\begin{example}[A shortage of eigenvectors]\label{ex:u23-defective}
$A=\begin{bmatrix} 3&1\\ 0&3\end{bmatrix}$ has $\det(A-\lambda I)=(3-\lambda)^{2}$, so
$\lambda_1=\lambda_2=3$. The eigenvectors fill the nullspace of
$A-3I=\left[\begin{smallmatrix} 0&1\\ 0&0\end{smallmatrix}\right]$, which has rank $1$, so
the nullspace is only the line through $(1,0)$. There is no second independent eigenvector.
Compare $3I$, which has the same eigenvalues $3,3$ but for which every nonzero vector is an
eigenvector. A repeated eigenvalue may or may not cost us eigenvectors; Section~\ref{sec:u24}
explains why this matters.
\end{example}

\subsection{Eigenvalues of powers, inverses and polynomials}

Once the eigenvectors of $A$ are known, many related matrices come for free.

\begin{proposition}[Functions of $A$]\label{prop:u23-functions}
Suppose $A\vb{x}=\lambda\vb{x}$. Then $A^{k}\vb{x}=\lambda^{k}\vb{x}$ for $k=1,2,\dots$, and
for any polynomial $p$, $p(A)\vb{x}=p(\lambda)\vb{x}$. If $A$ is invertible, then
$\lambda\neq0$ and $A^{-1}\vb{x}=\lambda^{-1}\vb{x}$.
\end{proposition}

\begin{proof}
$A^{2}\vb{x}=A(\lambda\vb{x})=\lambda A\vb{x}=\lambda^{2}\vb{x}$, and induction gives
$A^{k}\vb{x}=\lambda^{k}\vb{x}$. A polynomial $p(A)=c_0I+c_1A+\dots+c_mA^{m}$ then sends
$\vb{x}$ to $(c_0+c_1\lambda+\dots+c_m\lambda^{m})\vb{x}$. If $A$ is invertible it is not
singular, so $\lambda\neq0$; multiplying $A\vb{x}=\lambda\vb{x}$ by $A^{-1}$ and dividing by
$\lambda$ gives $A^{-1}\vb{x}=\lambda^{-1}\vb{x}$.
\end{proof}

\begin{example}[Eigenvalues of $A^{2}$ and $A^{-1}-I$ without computing them]\label{ex:u23-recitation}
Find the eigenvalues and eigenvectors of $A^{2}$ and of $A^{-1}-I$ for
\[
  A=\begin{bmatrix} 1&2&3\\ 0&1&-2\\ 0&1&4\end{bmatrix}.
\]

By Proposition~\ref{prop:u23-functions} it is enough to find those of $A$. Expand
$\det(A-\lambda I)$ down the first column, which has a single nonzero entry:
\[
  \det(A-\lambda I)=(1-\lambda)\begin{vmatrix} 1-\lambda & -2\\ 1 & 4-\lambda\end{vmatrix}
  =(1-\lambda)(\lambda^{2}-5\lambda+6)=(1-\lambda)(\lambda-2)(\lambda-3).
\]
The eigenvalues are $1,2,3$ (trace check: $1+1+4=6=1+2+3$). The eigenvectors:
\begin{itemize}
  \item $\lambda=1$: $A-I=\left[\begin{smallmatrix} 0&2&3\\ 0&0&-2\\ 0&1&3\end{smallmatrix}\right]$
  has a zero first column, so $x_1$ is free; the last two rows force $x_3=0$, then $x_2=0$.
  Eigenvector $(1,0,0)$.
  \item $\lambda=2$: $A-2I=\left[\begin{smallmatrix} -1&2&3\\ 0&-1&-2\\ 0&1&2\end{smallmatrix}\right]$.
  Row~2 gives $x_2=-2x_3$; row~1 gives $x_1=2x_2+3x_3=-x_3$. With $x_3=1$: $(-1,-2,1)$.
  \item $\lambda=3$: $A-3I=\left[\begin{smallmatrix} -2&2&3\\ 0&-2&-2\\ 0&1&1\end{smallmatrix}\right]$.
  Row~3 gives $x_2=-x_3$; row~1 gives $2x_1=2x_2+3x_3=x_3$. With $x_3=2$: $(1,-2,2)$.
\end{itemize}
These three vectors are eigenvectors of $A^{2}$ with eigenvalues $1,4,9$, and of $A^{-1}-I$
with eigenvalues $\tfrac11-1=0$, $\tfrac12-1=-\tfrac12$ and $\tfrac13-1=-\tfrac23$. We
never squared or inverted a $3\times3$ matrix.
\end{example}

\subsection{Exercises}

\begin{exercise}\label{exr:u23-B012}
A $3\times3$ matrix $B$ has eigenvalues $0$, $1$, $2$. Which of the following can be
determined from this information? Give them where possible.
(a) the rank of $B$; (b) $\det(B\T B)$; (c) the eigenvalues of $B\T B$;
(d) the eigenvalues of $(B^{2}+I)^{-1}$.
\end{exercise}
\begin{solution}
(a) The rank is $2$. Since $0$ is an eigenvalue, $B$ is singular and the rank is at most $2$.
Eigenvectors $\vb{x}_1,\vb{x}_2$ for $\lambda=1,2$ lie in the column space, because
$\vb{x}=B(\vb{x}/\lambda)$ when $\lambda\neq0$. They are independent: if
$\vb{x}_2=c\,\vb{x}_1$ then $2\vb{x}_2=B\vb{x}_2=c\,\vb{x}_1=\vb{x}_2$, impossible for
$\vb{x}_2\neq\vb{0}$. So the rank is at least $2$.

(b) $\det(B\T B)=\det B\T\det B=0\cdot0=0$, since $\det B=0\cdot1\cdot2$.

(c) Not determined. For $B=\operatorname{diag}(0,1,2)$, $B\T B=\operatorname{diag}(0,1,4)$
has eigenvalues $0,1,4$. For the triangular
$B=\left[\begin{smallmatrix} 0&1&0\\ 0&1&0\\ 0&0&2\end{smallmatrix}\right]$, which also has
eigenvalues $0,1,2$, we get
$B\T B=\left[\begin{smallmatrix} 0&0&0\\ 0&2&0\\ 0&0&4\end{smallmatrix}\right]$ with
eigenvalues $0,2,4$.

(d) By Proposition~\ref{prop:u23-functions}, an eigenvector of $B$ with eigenvalue
$\lambda$ is an eigenvector of $B^{2}+I$ with eigenvalue $\lambda^{2}+1\neq0$, hence of its
inverse with eigenvalue $1/(\lambda^{2}+1)$. The eigenvalues are $1$, $\tfrac12$, $\tfrac15$.
\end{solution}

\begin{exercise}\label{exr:u23-ABC}
Find the eigenvalues of
\[
  A=\begin{bmatrix} 1&2&3\\ 0&4&5\\ 0&0&6\end{bmatrix},\qquad
  B=\begin{bmatrix} 0&0&1\\ 0&2&0\\ 3&0&0\end{bmatrix},\qquad
  C=\begin{bmatrix} 2&2&2\\ 2&2&2\\ 2&2&2\end{bmatrix}.
\]
\end{exercise}
\begin{solution}
$A$ is triangular, so its eigenvalues are $1,4,6$
(Proposition~\ref{prop:u23-triangular}). For $B$, expand along the middle row:
\[
  \det(B-\lambda I)=(2-\lambda)\begin{vmatrix} -\lambda & 1\\ 3 & -\lambda\end{vmatrix}
  =(2-\lambda)(\lambda^{2}-3),
\]
so the eigenvalues are $2$, $\sqrt3$, $-\sqrt3$ (check: the sum is $2=\tr B$ and the product
$-6$ equals $\det B=-6$). $C$ has rank $1$, so its nullspace is a plane and $\lambda=0$ is an
eigenvalue twice; the trace $6$ gives the third eigenvalue $6$, with eigenvector $(1,1,1)$.
Directly, $\det(C-\lambda I)=-\lambda^{3}+6\lambda^{2}=-\lambda^{2}(\lambda-6)$.
\end{solution}

\begin{exercise}\label{exr:u23-idempotent}
Show that every eigenvalue of a projection matrix ($P^{2}=P$) is $0$ or $1$.
\end{exercise}
\begin{solution}
If $P\vb{x}=\lambda\vb{x}$ with $\vb{x}\neq\vb{0}$, then by
Proposition~\ref{prop:u23-functions} $\lambda^{2}\vb{x}=P^{2}\vb{x}=P\vb{x}=\lambda\vb{x}$,
so $(\lambda^{2}-\lambda)\vb{x}=\vb{0}$ and $\lambda(\lambda-1)=0$. Example~\ref{ex:u23-projection}
shows both values occurring: the plane gives $\lambda=1$, the normal line $\lambda=0$.
\end{solution}

\begin{exercise}\label{exr:u23-rotation}
Find the eigenvalues and eigenvectors of the rotation through the angle $\theta$,
\[
  R=\begin{bmatrix} \cos\theta & -\sin\theta\\ \sin\theta & \cos\theta\end{bmatrix},
\]
and check Theorem~\ref{thm:u23-trace}.
\end{exercise}
\begin{solution}
$\det(R-\lambda I)=(\cos\theta-\lambda)^{2}+\sin^{2}\theta=\lambda^{2}-2\lambda\cos\theta+1$.
The quadratic formula gives
$\lambda=\cos\theta\pm\sqrt{\cos^{2}\theta-1}=\cos\theta\pm i\sin\theta=e^{\pm i\theta}$.
For $\lambda=e^{i\theta}$, the vector $(1,-i)$ works:
$R(1,-i)=(\cos\theta+i\sin\theta,\;\sin\theta-i\cos\theta)=e^{i\theta}(1,-i)$; the conjugate
$(1,i)$ goes with $e^{-i\theta}$ (Proposition~\ref{prop:u23-conj}). The sum
$e^{i\theta}+e^{-i\theta}=2\cos\theta$ is the trace and the product $1$ is
$\det R=\cos^{2}\theta+\sin^{2}\theta$. For $\theta=90^{\circ}$ this is the matrix $Q$ above,
with eigenvalues $\pm i$. Both eigenvalues have $|\lambda|=1$: a rotation does not change
lengths.
\end{solution}
